\documentclass[11pt]{article}
\usepackage[utf8]{inputenc}
\usepackage[T1]{fontenc}
\usepackage[margin=1in]{geometry}
\usepackage{hyperref}
\usepackage{enumitem}
\title{STEM Challenge Lesson or FLEX Day as needed}
\author{Piter Garcia}
\date{March 20, 2026}
\begin{document}
\maketitle
\section*{Quick Scan}
\begin{itemize}[leftmargin=*]
\item Grade: Kindergarten
\item Workbook sessions: Session 6
\item Class blocks: Day 1 12:25-1:15 Justinger, Day 2 12:25-1:15 Kesys, Day 4 12:25-1:15 Pum
\item Reference file: stem-challenge-lesson-or-flex-day-as-needed.bib
\end{itemize}
\section*{School Planning Goals and Inclusion}
\begin{itemize}[leftmargin=*]
\item What is expected: I can rethink my ideas to improve my design | I can work with a partner | I can communicate my ideas | I can compromise or try my partners ideas | I can be supportive of my partner by collaborating with their design ideas. | I can sense my feelings and manage my feelings as I work
\item What we achieved: This lesson points to local transcript or evidence files in the workspace so the packet can be revised from what actually happened in class.
\item What needs improvement: stronger proof, cleaner assessment notes, and tighter lesson artifacts.
\item How to improve: add transcript proof, one saved artifact, and clearer success criteria.
\item Inclusion focus: use visual modeling, chunked directions, predictable routines, and flexible participation roles.
\end{itemize}
\section*{Official References}
\begin{itemize}[leftmargin=*]
\item \url{https://csteachers.org/k12standards/interactive/}

\end{itemize}
\end{document}
